\section{Discussion}

We have shown that US influenza surveillance data provide no
statistical evidence for inter-regional coupling between states. The
Cram\'er-Rao Lower Bound formalises this claim: for epidemics
entrained by shared seasonal forcing with similar initial conditions,
the Fisher Information for the coupling parameter approaches zero, so
the variance of any unbiased estimator is large. A
maximum-likelihood optimizer still converges and can return an
estimate that appears well-determined, but the CRLB-derived 95\%
confidence interval around it overlaps zero for nearly all 1081~state
pairs; the estimate is statistically indistinguishable from the null
hypothesis $\theta = 0$. The approximately uniform spread of
$\hat\theta$ across state pairs is symptomatic of this
uninformativeness rather than its definition: it is what one expects
when each pair's point estimate is set by the noise realization in
its 12~seasons rather than by any signal the data carry about
coupling.

\subsection{Why Connectivity Is Unidentifiable}

The mechanism is simple. Shared seasonal forcing \emph{entrains}
regional epidemics to a common trajectory---a form of resonance that
looks superficially like coupling-induced synchronization but requires
no coupling at all. When two entrained regions also start with
similar initial susceptibility, their incidence curves are nearly
indistinguishable regardless of $\theta$: a model with $\theta = 0$
(isolated) and a model with $\theta = 0.5$ (fully coupled) produce the
same trajectory, because any infection exchanged between regions is
offset by an equal flow in the opposite direction. The CRLB formalizes
this symmetry: $\partial\mu/\partial\theta \approx 0$, so the Fisher
Information vanishes.

\subsection{Recommendation for Endemic-Disease Preparedness Forecasting}

For endemic seasonal pathogens such as
influenza~\citep{viboud2006, tamerius2013}---pathogens that persist
at baseline levels year-round with annual epidemics driven by shared
environmental forcing---our two independent lines of evidence
converge on a single operational recommendation:
\textbf{inter-regional coupling is not a necessary component of
  preparedness forecasting models}. The CRLB analysis makes this a
statement about identifiability: no unbiased estimator of the
connectivity parameter can produce confidence intervals that
exclude zero, on any foreseeable amount of surveillance data, under
the entrainment regime that temperate US seasonal influenza actually
occupies. The seasonal-ARIMA benchmark makes it a statement about
forecast skill and simultaneously localises the mechanism: when a
bilateral cross-state covariate \emph{does} appear to help, it is a
noisy proxy for shared national-level flu-season severity, and that
signal is captured an order of magnitude more cleanly by a single
leave-one-out national-average covariate (median $7.5\%$ vs $0.7\%$
RMSE reduction; $97\%$ of states helped vs $64\%$ of pairs). The
relevant evidence is not the absolute size of the bilateral effect,
which would be uninformative in isolation, but its ratio to the
national-average baseline on the same data: the small bilateral
effect combined with the much larger pooled effect is precisely what
one would expect in the absence of directed coupling.

The implication for operational practice is direct. Seasonal-flu
forecasts---organised through the CDC FluSight
challenges~\citep{reich2019} and their analogues for other
pathogens~\citep{johansson2019}---inform hospital staffing, vaccine
distribution timing, antiviral stockpile management, and the timing
of public-health communication. They are not typically used to argue
for inter-state travel restrictions or quarantine. For this
preparedness purpose, forecasting models can safely replace
state-pair connectivity terms with a single national-severity
covariate (or equivalently, treat states as exchangeable conditional
on a shared severity signal), with no loss in forecast skill and a
substantial gain in interpretability and calibration. This also
offers a partial explanation for the well-documented difficulty of
epidemic forecasting~\citep{wilke2020, rosenkrantz2022}: any
forecast product that depends on an unidentified connectivity
parameter inherits that parameter's full uncertainty, so removing
the parameter tightens forecast characterisation rather than
loosening it.

Our result sits in productive tension with the recent empirical
comparison of~\citet{thivierge2025}, who asked precisely whether ILI
forecasting in US states benefits from including neighbour-state or
US-weighted-average ILI as covariates in quantile / linear / Poisson
autoregressions. They found modest but consistent improvements by
both RMSE and weighted interval score, and importantly observed that
neighbour-state and US-weighted-average covariates produce
\emph{similar} gains---an empirical hint, at their magnitude
resolution, that proximity is not the driving mechanism. Our results
sharpen that qualitative observation into a quantitative and
mechanistic claim. In our SARIMA framework the effects are not
similar: the leave-one-out national-average covariate improves
forecasts by an order of magnitude more than the typical bilateral
neighbour-state covariate ($7.5\%$ vs $0.7\%$ median RMSE
reduction; $97\%$ vs $64\%$ of forecasts helped). The CRLB analysis
supplies the mechanism: the connectivity parameter~$\theta$ is
structurally unidentifiable under entrainment, so the small
bilateral effect~\citet{thivierge2025} observe cannot be carrying
directional information about inter-regional transmission, and the
pooled-baseline dominance we report is what one expects when the
underlying signal is shared national-level seasonal severity rather
than directed coupling. Our paper therefore complements
\citet{thivierge2025} by providing both the information-theoretic
reason for what they observe and a refined decomposition of where
the usable signal actually lives. The related line of work that
pools weekly flu activity across regions to build forecast
benchmarks~\citep{viboud2003, yang2014, tsang2024} is consistent
with this picture: those benchmarks succeed because they exploit the
shared-severity signal, not inter-regional predictive coupling.

The contrast between temperate and tropical influenza dynamics maps
directly onto our framework. In temperate regions, shared winter
forcing entrains state-level epidemics, placing inference in the
high-CRLB regime. In tropical regions, weaker and less coherent
seasonal forcing produces asynchronous outbreak
patterns~\citep{tamerius2013}, which our simulations show would
provide the phase asymmetry needed for identifiability. Connectivity
inference from case data may therefore be more feasible in tropical
settings, or during the early invasion phase of a pandemic when
regional epidemics are naturally out of phase---precisely the regime
where early-outbreak interventions such as travel measures are
policy-relevant. The CRLB analysis makes this distinction precise.

\subsection{Beyond Influenza}

The identifiability check we apply here is not specific to
influenza or to inter-regional coupling. Contemporary spatial and
structured epidemic models routinely carry many
parameters---connectivity matrices, age-specific contact rates,
behavioural-response terms, waning-immunity kernels---each of which
requires that the surveillance data carry usable information about
its effect on observables. Because the CRLB bounds the variance of
every unbiased estimator and is computable from the model and a
hypothesised parameter value alone, the same design-stage check can
be applied parameter-by-parameter to any such model before any
inference pipeline is committed; parameters whose CRLB exceeds the
scale of their expected effect cannot be supported by the planned
data and should be fixed from external sources, or removed from the
model, rather than left as free degrees of freedom. The
mathematical structure underlying our particular
finding---oscillators driven by common periodic forcing---also
arises outside epidemiology, in neural circuits receiving shared
input~\citep{reid2019, pikovsky2001} and in metapopulations driven
by environmental cycles~\citep{bjornstad1999}; the same
entrainment-driven unidentifiability is plausible there, and the
same design-stage diagnostic applies.

\subsection{Limitations}

Several limitations warrant discussion. First, our model assumes a
single influenza strain with uniform transmission parameters across
all seasons and regions. In reality, $R_0$ varies substantially across
influenza subtypes~\citep{biggerstaff2014}. This variability is
partially absorbed by allowing the initial susceptible fraction
$S_i(0)$ to vary freely across seasons, but the seasonal forcing shape
is fixed.

Second, our observation model uses a Gaussian approximation to the
binomial. Real mortality data likely exhibit overdispersion; a negative
binomial model would increase the observation variance and inflate the
CRLB further. Since the mechanism driving unidentifiability---vanishing
sensitivity $\partial\mu/\partial\theta$ under symmetry---is
independent of the noise model, our Gaussian CRLB represents an
optimistic (smallest) bound. Any more realistic noise model would
reinforce the conclusion.

Third, we focused on a two-region model with symmetric coupling.
Extensions to asymmetric coupling or larger networks may reveal
additional identifiability challenges. Fourth, we compute excess
deaths by subtracting a simple pre-season
baseline~\citep{serfling1963}; more sophisticated baseline models
could improve the influenza-specific signal.
